% Problem record op_d3e38bc5b5b35af3. This TeX file is derived from
% database/problems_json/op_d3e38bc5b5b35af3.json by scripts/sync-tex.mjs; edit the JSON record
% and rerun the script rather than editing this file.
\section{Optimal precision dependence of diamond-norm channel tomography}
\label{sec:op_d3e38bc5b5b35af3}

\paragraph{Problem.}
Does tomography of an arbitrary $d$-dimensional quantum channel in diamond norm require and suffice with $\Theta(d^4/\varepsilon^2)$ channel uses?

For $d\geq2$, let $\Lambda:\mathcal L(\mathbb C^d)\to\mathcal L(\mathbb C^d)$ be an unknown completely positive, trace-preserving channel with no Kraus-rank promise. Let $Q_\diamond(d,\varepsilon)$ be the minimum worst-case number of ordinary channel uses needed to output $\widehat\Lambda$ such that

\begin{equation}
\Pr[\|\widehat\Lambda-\Lambda\|_\diamond\leq\varepsilon]\geq\frac23,
\qquad
\|\Phi\|_\diamond=\sup_\omega\| (\Phi\otimes\operatorname{id}_d)(\omega)\|_1.
\label{eq:d3e3-1}
\end{equation}

The supremum in Eq.~\eqref{eq:d3e3-1} is over states on $\mathbb C^d\otimes\mathbb C^d$. Adaptive inputs, ancillas, quantum memory, and collective measurements are allowed, but no purification of the channel environment is supplied. Is $Q_\diamond(d,\varepsilon)=\Theta(d^4/\varepsilon^2)$ uniformly in $d$ and sufficiently small $\varepsilon$?

\subsection*{Status}

Solved

\subsection*{Source}

This is an editor-formulated unrestricted-rank specialization of the joint dimension--accuracy question discussed by Mele and Bittel \sourcecite{ref:d3e3-mele25}{Mele25}. Subsequent lower bounds settle this specialization; the original question is retained here as a solved problem.

\subsection*{Progress}

\begin{itemize}
  \item General channels of input dimension $d_1$, output dimension $d_2$, and Kraus rank at most $r$ admit diamond-norm tomography with $O(rd_1d_2/\varepsilon^2)$ ordinary channel queries \sourcecite{ref:d3e3-mele25}{Mele25}, \sourcecite{ref:d3e3-chen25}{Chen25}. Setting $d_1=d_2=d$ and $r=d^2$ gives the required upper bound.

  \item Chen, Zhang, and Yu prove the matching lower bound $\Omega(rd_1d_2/\varepsilon^2)$ when $rd_2\geq2d_1$ (Theorem~1.1 and Corollary~1.2 of the January 2026 preprint) \sourcecite{ref:d3e3-czy26}{CZY26}. The proof uses sequential channel testers, so it permits adaptive coherent queries and quantum memory; it is not restricted to tomography of independently prepared Choi states.

  \item The subsequent joint preprint of Chen, Girardi, Oufkir, Yu, and Zhang subsumes that result. Corollary~1.3 gives $\Theta(rd^2/\varepsilon^{\min\{r,2\}})$ for equal input and output dimensions; Theorems~3.4 and~6.2 supply the upper and lower bounds, with the query model defined in Section~2.3. Substituting $r=d^2$ yields \begin{equation}Q_\diamond(d,\varepsilon)=\Theta(d^4/\varepsilon^2).\label{eq:d3e3-optimal}\end{equation} The rate is uniform in the dimension and sufficiently small error \sourcecite{ref:d3e3-cgoyz26}{CGOYZ26}. These are preprint results.
\end{itemize}

\subsection*{References}

\begin{enumerate}[label={},leftmargin=4.8em,labelsep=0.5em]
  \footnotesize
  \item[\textup{[Mele25]}]\label{ref:d3e3-mele25}
  A. A. Mele and L. Bittel, "Optimal learning of quantum channels in diamond distance," arXiv preprint (2025), version 3, 15 June 2026. \href{https://arxiv.org/abs/2512.10214}{arXiv:2512.10214}.

  \item[\textup{[Chen25]}]\label{ref:d3e3-chen25}
  K. Chen, N. Yu, and Z. Zhang, "Quantum channel tomography and estimation by local test," arXiv preprint (2025), version 2, 2 February 2026. \href{https://arxiv.org/abs/2512.13614}{arXiv:2512.13614}.

  \item[\textup{[CZY26]}]\label{ref:d3e3-czy26}
  K. Chen, Z. Zhang, and N. Yu, ``Optimal lower bound for quantum channel tomography in away-from-boundary regime,'' arXiv preprint, version 1 (2026). \href{https://arxiv.org/abs/2601.10683v1}{arXiv:2601.10683v1}.

  \item[\textup{[CGOYZ26]}]\label{ref:d3e3-cgoyz26}
  K. Chen, F. Girardi, A. Oufkir, N. Yu, and Z. Zhang, ``Quantum channel tomography: optimal bounds and a Heisenberg-to-classical phase transition,'' arXiv preprint, version 2, 9 July 2026. \href{https://arxiv.org/abs/2604.17369v2}{arXiv:2604.17369v2}.
\end{enumerate}

\subsection*{Comment}

The product dependence is established, not merely the two one-parameter scalings separately. The lower bound covers the ordinary black-box model in the statement, including adaptive queries, and the upper bound needs no access to a purification of the channel environment. The dimension is taken to be $d\geq2$; at $d=1$ there is a unique channel and zero queries suffice. The explicit large-dimension cutoff in Theorem~6.2 does not change the asymptotic claim: for the finitely many smaller nontrivial dimensions, preparation channels reduce the task to state tomography and give a uniform lower bound after decreasing its constant. The cited resolution predates this weekly audit and corrects an omission from the previous progress summary.

\subsection*{Field}

Quantum metrology

\subsection*{Topic}

Quantum estimation; One-shot and finite-blocklength bounds

\subsection*{ID}

\texttt{op\_d3e38bc5b5b35af3}
